\section{Related work}\label{sec:related}

\paragraph{Low-light video enhancement.}
Video restorers combine spatial reconstruction with temporal redundancy through
decomposition~\citep{stcd2025}, spatial--temporal compatible
learning~\citep{tcenet}, recurrence~\citep{retinexmcnet}, state-space
models~\citep{bvi_mamba}, or diffusion~\citep{bvi_rlv,dl_diff}.
Image methods provide strong frame-wise alternatives, including
Retinexformer~\citep{retinexformer}, HVI-CIDNet~\citep{hvi_cidnet},
DarkIR~\citep{darkir}, and zero-reference methods~\citep{zerodce2020,sci2022}.
BVI-RLV~\citep{bvi_rlv} makes sensor and scene transfer explicit.
Our main comparison holds the restoration architecture and training data fixed
and changes its conditioning; a separately retrained Retinexformer provides
context rather than a claim of universal superiority.

\paragraph{Frozen representations for restoration.}
Pretrained features can judge reconstruction quality through perceptual
losses~\citep{johnson_perceptual,lpips}, or condition the restorer directly.
DA-CLIP~\citep{daclip} uses adapted vision--language embeddings, and pretrained
features also guide restoration networks~\citep{xu2024priors} and diffusion
restorers~\citep{stablesr}. We study the latter role with a frozen video
encoder: V-JEPA~\citep{vjepa,vjepa2}, learned by predicting masked
spatio-temporal representations. Its training objective motivates a temporal
prior, but does not establish that a downstream restorer uses temporal
information. We measure that distinction by intervening on the tokens and on
the training objective.

\paragraph{Representation versus downstream use.}
A representation may contain information that a particular readout does not
use. Our experiments therefore separate encoder sensitivity, restoration
quality, and the sensitivity of a trained restorer to token interventions.
Likewise, low rank after training is a property of that model under a specified
intervention; it does not by itself establish a smaller trainable model, a
cheaper encoder, or equivalence between two conditioning distributions.
